\documentclass[10pt,a4paper]{amsart}
\usepackage[margin=19mm]{geometry}
\usepackage{amsmath,amssymb,mathtools}
\usepackage[scr=rsfs,cal=euler]{mathalfa}
\usepackage{xfrac}
\usepackage[unicode,hidelinks,bookmarks=false]{hyperref}
\newcommand{\ie}{i.e.\ }
\newcommand{\cf}{cf.\ }
\newcommand{\resp}{respectively\ }
\newcommand{\Subsection}{\subsection}
\providecommand{\nobreakdash}{\nobreak\hbox{-}}
\setlength{\emergencystretch}{2em}
\multlinegap0pt
\usepackage{amssymb}
\usepackage{enumerate}
\usepackage{float}
\newtheorem{proposition}{Proposition}
\def\bR{\mathbb{R}}
\def\cH{\mathcal{H}}
\title{Failure of zero extension in parabolic Sobolev spaces}
\author{Jongkeun Choi, Doyoon Kim,, Kwan Woo}
\date{Source: arXiv:2606.22059v1 (CC BY 4.0)}
\begin{document}
\maketitle

\noindent\emph{Source and licence.} Failure of zero extension in parabolic Sobolev spaces. Version arXiv:2606.22059v1 (CC BY 4.0), distributed by arXiv under the Creative Commons Attribution 4.0 International licence, http://creativecommons.org/licenses/by/4.0/, as recorded in the arXiv licence metadata for this exact version and shown on the abstract page https://arxiv.org/abs/2606.22059v1. Full licence notice: \texttt{licenses/arxiv-2606.22059-CC-BY-4.0.txt}.\par\smallskip

\noindent\textit{Editorial scope.} Counted unit: the article's proposition prop\_inf (source0.tex lines 993--998). The article's complete original proof of it is delivered with it (source0.tex lines 999--1064, one \texttt{proof} environment of 66 lines). No mathematics is written, repaired or re-derived: the statement and the proof are the article's own lines, in the article's own notation, and no retained line is substituted or re-flowed.  No further source-local context is needed: every cross-reference of every retained line resolves inside the two delivered spans.  Nothing was omitted from any retained span, so the package carries no omission record.  No retained line contains a citation, so the wrapper carries no bibliography and no bibliography entry.  No retained line contains a figure environment, an image-inclusion command or drawing code, so no image is delivered and the package declares no asset dependency.  Editorial formatting only, in addition: an AMS \texttt{amsart} preamble that restates the article's own declarations and macros used by the retained lines, namely the environments \texttt{proposition} on separate counters, together with the standard packages those lines need.  The retained environments print in this document as Proposition 1 in order of appearance, whereas the article prints the same items with the numbering of its own full document.  The complete native source of the article as distributed in the arXiv e-print for this exact version is bundled unchanged as \texttt{sources/203319/source0.tex}.
\begin{proposition} \label{prop_inf}
Let $T \in (0, \infty]$. For any $u \in \mathring{\cH}^1_\infty((0,T)\times \bR^d_+)$, its zero extension $\bar{u}$ with respect to $x$ satisfies
$$
\bar{u} \in \cH^1_\infty((0,T)\times \bR^d).
$$
\end{proposition}

\begin{proof}
Since $u \in \mathring{\cH}^1_\infty((0,T)\times \bR^d_+)$, we have $u, Du \in L_\infty((0,T)\times \bR^d_+)$, and there exist $g_1, \ldots, g_d, f \in L_\infty((0,T)\times \bR^d_+)$ such that
\begin{equation}		\label{260616_eq1}
u_t = D_i g_i + f \quad \text{in }\, (0,T)\times \bR^d_+.
\end{equation}
For the zero extension $\bar{u}$, it is clear that $\bar{u}, D\bar{u} \in L_\infty((0,T)\times \bR^d)$.
To show that $\bar{u}_t$ can be represented in divergence form, let $\varphi \in C^\infty_0((0,T)\times \bR^d)$ be an arbitrary test function.
We consider the functional
$$
I = I(\varphi) := -\int_0^T \int_{\bR^d} \bar{u} \varphi_t \,dx\,dt = -\int_0^T \int_{\bR^d_+} u \varphi_t \,dx\,dt.
$$

For $\varepsilon > 0$, we introduce a cut-off function $\zeta_\varepsilon(x_1)$ defined by
\[
\zeta_\varepsilon(x_1) = 
\begin{cases}
    0 & \text{for } x_1 \le 0, \\[4pt]
    x_1/\varepsilon & \text{for } 0 < x_1 < \varepsilon, \\[4pt]
    1 & \text{for } x_1 \ge \varepsilon.
\end{cases}
\]
Then by testing \eqref{260616_eq1} with $\zeta_{\varepsilon}\varphi$, we have
$$
-\int_0^T \int_{\bR^d_+} u (\zeta_\varepsilon \varphi)_t \,dx\,dt = \int_0^T \int_{\bR^d_+} - g_i D_i(\zeta_\varepsilon \varphi) + f \zeta_\varepsilon \varphi  \,dx\,dt.
$$
Taking the limit as $\varepsilon \to 0$, the left-hand side converges to $I$. 
On the right-hand side, since $D_1 \zeta_\varepsilon = \varepsilon^{-1} I_{(0, \varepsilon)}(x_1)$, we obtain
\begin{equation}
    \label{eq:p_inf}
I = \int_0^T \int_{\bR^d_+} - g_i D_i \varphi + f \varphi  \,dx\,dt + \lim_{\varepsilon \to 0^+} \frac{1}{\varepsilon} \int_0^T \int_0^\varepsilon \int_{\bR^{d-1}} g_1 \varphi \,dx'\,dx_1\,dt.
\end{equation}
Note that
$$
|\varphi(t, x_1, x')| = \bigg| \int_{-\infty}^{x_1} D_1 \varphi(t, s, x') \,ds \bigg| \le \int_{-\infty}^\infty |D_1 \varphi(t, s, x')| \,ds.
$$
Using this, we estimate the last term in \eqref{eq:p_inf} as follows:
$$
\begin{aligned}
&\bigg| \frac{1}{\varepsilon} \int_0^T \int_0^\varepsilon \int_{\bR^{d-1}} g_1 \varphi \,dx'\,dx_1\,dt \bigg| \\
&\le \|g_1\|_{L_\infty((0,T)\times \bR^d_+)} \frac{1}{\varepsilon} \int_0^\varepsilon \bigg( \int_0^T \int_{\bR^{d-1}} \int_{-\infty}^\infty |D_1 \varphi(t,s, x')| \,ds \,dx'\,dt \bigg) dx_1 \\
&= \|g_1\|_{L_\infty((0,T)\times \bR^d_+)} \int_0^T \int_{\bR^d} |D_1 \varphi(t, x)| \,dx\,dt.
\end{aligned}
$$
Therefore, 
$$
|I| \le N \int_0^T \int_{\bR^d} \bigg( |\varphi| + \sum_{i=1}^d |D_i \varphi| \bigg) \,dx\,dt,
$$
where
\[
N=N\left(\|g_i\|_{L_\infty((0,T)\times \bR^d_+)}, \|f\|_{L_\infty((0,T)\times \bR^d_+)}\right) > 0.
\]
This inequality implies that the linear mapping
\[
(\varphi, D_1 \varphi, \dots, D_d \varphi) \mapsto I
\]
is bounded with respect to the $L_1$-norm.
By the Hahn-Banach theorem and the duality relation $(L_1)^* = L_\infty$, there exist functions $\bar{F}, \bar{G}_1, \dots, \bar{G}_d \in L_\infty((0,T)\times \bR^d)$ such that
$$
I = \int_0^T \int_{\bR^d} -\bar{G}_i D_i \varphi +  \bar{F} \varphi \,dx\,dt.
$$
This concludes that
\[
\bar{u}_t = D_i \bar{G}_i + \bar{F} \quad \textrm{ in }\,\, (0,T)\times \bR^d.
\]
Hence, $\bar{u} \in \cH^1_\infty((0,T)\times \bR^d)$.
\end{proof}

\end{document}
